%% Contextual predicate logic: finite variables and finite hypotheses differ.

\subsection{Contextual first-order syntax and its interpretation}
\label{sec:p10-first-order-categorical-semantics}

\begin{hbdefinition}{def:p10-contextual-first-order}{The contextual first-order language and calculus}
Fix a many-sorted signature with a set of sorts, function symbols
\(F:\sigma_1,\ldots,\sigma_n\to\tau\), and relation symbols
\(R\subseteq\sigma_1,\ldots,\sigma_n\), each of finite arity,
including zero. Fix a countably infinite supply of variable names;
their sorts are assigned by declarations. A variable context \(\Gamma\) is a finite list of
distinct variables with declared sorts. Write \(M:\sigma\,[\Gamma]\)
for a well-sorted term: a declared variable has its declared sort,
and \(F(M_1,\ldots,M_n)\) has sort \(\tau\) when its arguments have
the specified sorts in that same context.

Following Pitts, use the formula constructors
\[
 R(\mathbf M),\quad (M=_\sigma N),\quad\mathsf{false},\quad
 \varphi\mathbin{\&}\psi,\quad\varphi\vee\psi,\quad
 \varphi\Rightarrow\psi,\quad
 \forall_\sigma(\varphi),\quad\exists_\sigma(\varphi).
\]
Atomic arguments must have their declared sorts, and equality requires
two terms of the same sort. The propositional constructors preserve the
variable context. In the last two constructors, \(\varphi\) is a
unary formula abstraction: if \(e\) is a formula in
\(\Gamma,x:\sigma\), with \(x\notin\operatorname{var}(\Gamma)\),
write \(\varphi=(x)e\) and \(\varphi(M)=e[M/x]\).
Then \(\varphi(x)\) is its body in \(\Gamma,x:\sigma\), while
\(\forall_\sigma(\varphi)\) and \(\exists_\sigma(\varphi)\) are
formulas in \(\Gamma\). This is Pitts's metalevel abstraction
notation \((x)e\), introduced in \cite[\S3, p.~13]{Pitts1995Draft};
the abstraction is notation for binding and substitution, not an
additional object-language sort or connective.
Put \(\neg\varphi=(\varphi\Rightarrow\mathsf{false})\) and
\(\mathsf{true}=(\mathsf{false}\Rightarrow\mathsf{false})\).
Formulas are taken up to sort-preserving bound alpha-renaming.
Every substitution is simultaneous and capture-avoiding; first freshen
bound variables outside the finitely many contexts and terms involved.

The judgments are term equations \(M=N:\sigma\,[\Gamma]\) and sequents
\[
                         \Phi\vdash\psi\,[\Gamma].
\]
Here \(\Phi\) is a finite list of assumptions, and every assumption
and the conclusion is a formula in \(\Gamma\).
The variable context and the hypothesis list have different roles.
In particular, the judgment retains unused variable declarations.
A theory \(\mathcal T\) is any set of well-formed term-equation and
sequent axioms. A derivation is a finite tree using those axioms and
Sheet~\ref{sheet:contextual-first-order}.
The sheet records the rule package of Pitts's Figures~1--3, 5,
and~7 with variable contexts and side conditions explicit.
Its displayed witness typing \(M:\sigma\,[\Gamma]\) in
\(\exists\)-Intr states the source's implicit well-sorted
instantiation condition; Figure~5 omits that premise from the printed
fraction. The equivalent adjoint presentations in Figures~4, 6,
and~8 are displayed separately in
Section~\ref{sec:p10-adjoint-first-order-calculus}, with their
syntactic interderivability proofs; they are not additional
primitive rules of this presentation.
\index{first-order logic!contextual}\index{variable context}\index{hypothesis list}
\end{hbdefinition}

\showcalculus{contextual-first-order}

The term-equation judgment and the formula \((M=_\sigma N)\) remain
distinct. There is no rule reflecting a proved equality formula into
an equation of term arrows. This is the distinction made in
\cite[\S5.6, Remark~5.6.7, p.~59]{Pitts1995Draft}.

\begin{hbdefinition}{def:p10-contextual-doctrine-interpretation}{Interpreting a contextual first-order signature}
Let \(P:\mathcal C^{\mathrm{op}}\to\mathsf{Heyt}\) be a first-order
hyperdoctrine with equality as in
Definitions~\ref{def:p10-first-order-hyperdoctrine} and
\ref{def:p10-doctrine-equality}.
Interpret each sort \(\sigma\) by an object \(A_\sigma\).
For a finite sort list use the chosen left-associated product starting
with \(1\); likewise put
\[
 I_{()}=1,\qquad I_{\Gamma,x:\sigma}=I_\Gamma\times A_\sigma.
\]
A unary variable context therefore denotes \(1\times A_\sigma\).
Interpret an \(n\)-ary function symbol by an arrow from its argument
product to its result object, and an \(n\)-ary relation symbol by a
predicate in the fibre over that argument product.
At arity zero these are respectively an arrow from \(1\) and a
predicate in \(P(1)\).

Terms denote arrows from \(I_\Gamma\): variables are projections and
\[
 \llbracket F(\mathbf M)\rrbracket_\Gamma
       =F^{\mathcal I}\circ\langle\llbracket M_i\rrbracket_\Gamma\rangle_i.
\]
The tuple into a leading-terminal product includes its unique terminal
component. In particular the empty tuple is \(!:I_\Gamma\to1\).
Formulas denote predicates in \(P(I_\Gamma)\):
\[
\begin{aligned}
 \llbracket R(\mathbf M)\rrbracket_\Gamma
   &=\langle\llbracket M_i\rrbracket_\Gamma\rangle_i^*R^{\mathcal I},\\
 \llbracket M=_\sigma N\rrbracket_\Gamma
   &=\langle\llbracket M\rrbracket_\Gamma,
                         \llbracket N\rrbracket_\Gamma\rangle^*\delta_{A_\sigma},\\
 \llbracket\mathsf{false}\rrbracket_\Gamma&=\bot_{I_\Gamma},\\
 \llbracket\varphi\mathbin{\&}\psi\rrbracket_\Gamma
   &=\llbracket\varphi\rrbracket_\Gamma\wedge\llbracket\psi\rrbracket_\Gamma,\\
 \llbracket\varphi\vee\psi\rrbracket_\Gamma
   &=\llbracket\varphi\rrbracket_\Gamma\vee\llbracket\psi\rrbracket_\Gamma,\\
 \llbracket\varphi\Rightarrow\psi\rrbracket_\Gamma
   &=\llbracket\varphi\rrbracket_\Gamma
                        \Rightarrow\llbracket\psi\rrbracket_\Gamma,\\
 \llbracket\forall_\sigma(\varphi)\rrbracket_\Gamma
   &=\forall_{\pi_{I_\Gamma,A_\sigma}}
                      \llbracket\varphi(x)\rrbracket_{\Gamma,x:\sigma},\\
 \llbracket\exists_\sigma(\varphi)\rrbracket_\Gamma
   &=\exists_{\pi_{I_\Gamma,A_\sigma}}
                      \llbracket\varphi(x)\rrbracket_{\Gamma,x:\sigma}.
\end{aligned}
\]
The implication on the right acts in a Heyting fibre; the one on the
left is the formula constructor.
A term equation is satisfied when its arrows are equal.
A sequent is satisfied when
\[
          \bigwedge_{\phi\in\Phi}\llbracket\phi\rrbracket_\Gamma
                                  \leq\llbracket\psi\rrbracket_\Gamma,
\]
with empty meet \(\top_{I_\Gamma}\).
\end{hbdefinition}

\begin{hbtheorem}{thm:p10-doctrine-substitution}{proved=annex:p10.quant-fol.1}{Terms, formulas, and contextual substitution}
Let \(\mathbf s:\Delta\to\Gamma\) be a tuple of well-sorted terms
in \(\Delta\), one for each variable of \(\Gamma\), and let
\(s:I_\Delta\to I_\Gamma\) be its interpreted tuple arrow.
Substitution preserves well-formedness, respects bound alpha-renaming,
and satisfies
\[
 \llbracket M[\mathbf s]\rrbracket_\Delta
       =\llbracket M\rrbracket_\Gamma\circ s,\qquad
 \llbracket\varphi[\mathbf s]\rrbracket_\Delta
       =s^*\llbracket\varphi\rrbracket_\Gamma.
\]
Adding unused variable declarations is reindexing along their product
projection; consistent renaming and permutation use the corresponding
variable-tuple isomorphisms.
\end{hbtheorem}

\begin{hbtheorem}{thm:p10-doctrine-soundness}{proved=annex:p10.quant-fol.2}{Soundness of contextual intuitionistic first-order logic}
Every interpretation in the stated doctrine satisfying every axiom of
\(\mathcal T\) satisfies every term equation and every contextual
sequent derived from \(\mathcal T\) on the displayed sheet.
The set of theory axioms may be arbitrary; each derivation uses only
finitely many. Sort objects need not be inhabited.
\end{hbtheorem}

The interpretations and stability conditions are
\cite[\S\S5.2--5.6, pp.~40--59, especially Lemma~5.2.6,
Theorems~5.3.2 and~5.4.2]{Pitts1995Draft}.
The proof below treats the displayed calculus directly, without
identifying it with a context-erasing presentation of first-order logic.

\begin{hbtheorem}{thm:p10-doctrine-tarski}{proved=annex:p10.quant-fol.3}{The powerset interpretation and nonempty sorts}
In the powerset doctrine over \(\mathsf{Set}\), the interpretation of
each formula in \(\Gamma\) is exactly the set of finite assignments
to \(\Gamma\) satisfying its usual many-sorted Tarskian clauses,
transported through the canonical product-to-tuple bijection.
This assertion allows empty sort carriers.
Term-equation and sequent satisfaction agree with universal validity
over these finite assignments.

When every interpreted sort is nonempty, these are the ordinary
nonempty-sort Tarskian semantics. An unused variable declaration can
then be removed from the semantic context of either kind of judgment.
For predicate sequents this removal is also valid in any doctrine
for inhabited sort objects, by Theorem~\ref{thm:p10-doctrine-inhabited}.
That abstract statement reflects predicate order; it does not assert
reflection of equality between base arrows.
Without this hypothesis, neither context removal nor validity of
the closed formula \(\exists_\sigma((x)(x=_\sigma x))\) is supplied.
\end{hbtheorem}

For an empty sort in Set, the context \((x:\sigma)\) has empty
interpretation, so even \(\vdash\mathsf{false}\,[x:\sigma]\) is satisfied
in that interpretation. The empty context denotes a singleton, where
\(\vdash\mathsf{false}\,[()]\) is not satisfied.
This is why the retained context is part of the semantic judgment.
The powerset comparison is Example~5.5.3, pp.~53--54, together with
Example~5.6.3, p.~56, of \textcite{Pitts1995Draft}; no completeness
claim for the current handbook LJ calculus is made here.
